\subsection{Negative-line-bundle quotients}
\label{sec:elliptic-line-endpoint-examples}
\index{Elliptic fibration!negative-line-bundle quotient}

The simplest examples of the curvature criterion come from negative line bundles. Dividing the complement of the zero section of a negative line bundle over $\mathbb P^1\times\mathbb P^1$ by a dilation produces a principal elliptic bundle, and the logarithm of the fiber norm gives a connection form whose curvature is a multiple of a K\"ahler form on the base.

Take $B=\mathbb P^1\times\mathbb P^1$, let $\omega_1,\omega_2$ be the pullbacks of the Fubini--Study forms of the two factors, normalized to have integral one, and put
\[
 \begin{gathered}
 \Omega=a\omega_1+b\omega_2,\quad a,b\in\mathbb N_{>0},\qquad
 L=\mathcal O(-a,-b),\\
 X=L^\times/\langle v\mapsto\rho v\rangle,\quad0<\rho<1,
 \end{gathered}
\]
where $L^\times$ is the total space of $L$ with the zero section removed. Equip $L$ with the Hermitian norm induced by the standard norms on the tautological bundles of the two factors. Let $r=\log\|v\|^2$ and set $\eta=\partial r$.

\begin{proposition}\label{prop:elliptic-line-endpoint-example}
The quotient $X$ is a compact complex threefold, and the projection to $B$ is a principal elliptic fibration. For $g=\pi^*\Omega+s\mathrm i\eta\wedge\bar\eta$, $s>0$,
\begin{equation}\label{eq:elliptic-line-lower-mass}
 v_-(g)=24\pi s ab\log(1/\rho).
\end{equation}
Every Hermitian form on $X$ satisfies the endpoint bound of \cref{thm:elliptic-curvature-endpoint} over its full set of normalized Gauduchon currents. The manifold has $H^2(X,\R)\ne0$ and admits no balanced or pluriclosed metric. In particular, it is not diffeomorphic to a Hopf threefold $(\C^3\setminus\{0\})/\langle z\mapsto\lambda z\rangle$, $0<|\lambda|<1$.
\end{proposition}
\begin{proof}
In product coordinates, the squared norm of the standard local frame of $L$ is $(1+|z_1|^2)^a(1+|z_2|^2)^b$. Therefore
\[
 \ddc r=\pi^*\Omega,\qquad \mathrm d\eta=2\pi\mathrm i\pi^*\Omega.
\]
Dilation changes $r$ by a constant, so $\eta$ descends to $X$. The function $r$ changes by the nonzero constant $\log\rho^2$ under the generator, so the action is free and properly discontinuous. The product of the compact unit circle bundle with a fundamental interval for $r$ maps onto $X$, which is therefore compact. Each fiber is $\C^*/\rho^{\mathbb Z}$. Its logarithmic coordinate has periods $\log\rho$ and $2\pi\mathrm i$, and the curvature $F=2\pi\mathrm i\Omega$ has real part $0$ and imaginary part $2\pi\Omega$; hence
\[
 A_E=4\pi\log(1/\rho),\qquad Q=4\pi^2\Omega^2>0,
 \qquad\int_B\Omega^2=2ab.
\]
\Cref{thm:elliptic-curvature-endpoint} now gives \eqref{eq:elliptic-line-lower-mass} and the endpoint bound, and excludes pluriclosed metrics.

To prove $H^2(X,\R)\ne0$, put $A=\eta+\bar\eta$ and $B_1=\mathrm i(\eta-\bar\eta)$. Then $\mathrm dA=0$ and $\mathrm dB_1=-4\pi\pi^*\Omega$. Suppose that $\pi^*(\omega_1-\omega_2)=\mathrm d\zeta$ for a real one-form $\zeta$. Averaging $\zeta$ over the compact group of elliptic translations, we may assume $\zeta$ invariant, so that
\[
 \zeta=\pi^*\xi+(\pi^*f)A+(\pi^*j)B_1
\]
for a one-form $\xi$ and functions $f,j$ on $B$. The components of $\mathrm d\zeta$ with one vertical and one horizontal direction must vanish, which gives $\mathrm df=\mathrm dj=0$. Thus
\[
 \omega_1-\omega_2=\mathrm d\xi-4\pi j\Omega.
\]
Integrating over the two rulings gives $1=-4\pi ja$ and $-1=-4\pi jb$, which is impossible because $a,b>0$. A Hopf threefold as above is diffeomorphic to $S^5\times S^1$, whose second real cohomology vanishes.

Finally, $\pi^*\Omega=-\mathrm dB_1/(4\pi)$ is a nonzero semipositive exact form. Its integral against the closed positive square of a balanced metric would vanish by Stokes' theorem and be strictly positive by positivity. Hence no balanced metric exists.
\end{proof}
